\scanpage{25}
% archive sheet number: 24

\textcircled{6} Ce qui se passe si on se restreint au cas où $k'$ est un corps [on est tenté de le faire au moins si $k$ en est un]. \deleted{\unsure{la correspondance}} Soit $V'_S=V_S/\text{corps sur }k$

\[
V^{\prime}_{\deleted{\mathrm{corps}},S}=V^{\prime}_{\deleted{\mathrm{corps}},J}=V^{\prime}_{\deleted{\mathrm{corps}},\widetilde J}
\]

Donc $J\mapsto V^{\prime}_{\deleted{\unsure{\mathrm{corps}}}J}$ plus injective. Mais sa restriction à l'ens des idéaux \deleted{\unsure{J}} $J$ tels que $J=\widetilde J$ l'est, car on a

\[
\widetilde J=\{f\in P_I\mid f(x)=0\quad\forall\,x\in V'_J(k'),\ k'/k\text{ un corps}\}
\]

[Cela \deleted{signifie} équivaut à l'énoncé suivant d'AC

\uline{Prop} $(\forall\,\underset{\text{idéal}}J\subset P_I)$ $\widetilde J$ est l'intersection des noyaux des hom de $P_I/J$ dans des corps

\[
\Downarrow
\]

$\forall$ anneau $A$, l'intersection des hom de $A$ dans des corps [i.e. l'intersection des idéaux premiers) est l'ens des éléments nilpotents. ] Donc

\[
\widetilde J\subset\widetilde J'\ \unsure{\longleftrightarrow|\Longleftrightarrow}\ V(J')\unsure{\subset|=}\,V(J)
\]

\unsure{M$^{\mathrm s}$|Mêmes} conclusions \inserted{pour $I$ fini} si on se borne à tte sous-catégorie \inserted{$C$} de la catégorie des $k$-algèbres réduites, pourvu toutefois que \unsure{$C$|$\mathcal C$|$k$} soit assez grande pour que l'analogue de la prop. précédente soit valable

\textcircled{$6'$} (par ex. $C$ contient les extensions de corps,] -- \deleted{\unsure{ou si}} si $I$ est fini, s'il contient les ext. finies des corps résiduels de $k$, -- p. ex. ou \unsure{aust. : $k$ si|aust : $k$ si|aura : $k$ si} $k$ est un corps alg. clos \ldots)

\margin{\uline{Nullstellensatz}}

\margin{Dans une alg. de t.f. sur corps $k$ tt idéal premier (réduit) est intersection d'idéaux maximaux, et le quot\unsure{ient} par un \ldots{} est \ldots}

\textcircled{7} Importance des \deleted{\unsure{points}} éléments nilpotents.

\[
\text{\textcircled{8}}\quad A\simeq\mathrm{Hom}_{k\text{-}\mathrm{alg}}(k[T],A)\simeq\mathrm{Hom}(V_A,\mathbb E_k^1)
\]

% slanted
\margin{définir d'abord $f(x)$ pr $f\in A$, $x\in V_A(k')=X(k')$ \quad $(fg)(x)=f(x)g(x)$, \unsure{$1(x)=1_{k'}$|$(1)(x)=1_{k'}$}, $(f+g)(x)=f(x)+g(x)$}

[si $A=P_I/J$, $f\in A$ se relève en $g\in P_I$, alors \deleted{\unsure{$\longleftrightarrow$}}

\diagram{\begin{tikzcd}[ampersand replacement=\&,row sep=small]
V_A(k')\arrow[r,"f_{k'}"]\&k'\&\text{est induit par}\\
\mathbb E^I(k')=k^{\prime I}\arrow[ur,"g_{k'}"']\&\&\text{défini par }g\,]
\end{tikzcd}}

\margin{$A$ \uline{l'anneau des fonctions sur l'espace algébrique affine} $X$.}

La structure d'anneau de $A$ provient de la structure d'anneau sur le foncteur $\mathbb E_k^1$.
